\section{Scratch: KL Chain Bounds and Summability}

This section records the pen-and-paper proof obligations behind the Rocq lemmas
\texttt{iterated\_kl\_chain\_bound} and
\texttt{summable\_kl\_from\_coordinate\_bounded\_kl}.  It is scratch text: the
goal is to make the shared analytic invariant visible before committing to the
engineering details.  The two results should be developed in parallel, since
they are value-level and summability-level faces of the same chain-rule
argument.

\paragraph{Setup.}
Let $P,Q$ be probability distributions on $A^n$.  For each coordinate $i<n$
and prefix $a\in A^i$, write $P_i(\cdot\mid a)$ and $Q_i(\cdot\mid a)$ for the
conditional distributions of coordinate $i$ given prefix $a$, using the zero
distribution when the conditioning event has zero mass.  Assume $P\ll Q$ and
\[
  \KL{P_i(\cdot \mid a)}{Q_i(\cdot \mid a)}
  \leq \varepsilon_i
  \qquad\text{for all } i,a.
\]
When proving summability, also assume these conditional KL divergences are
finite.

\begin{lemma}[Iterated KL chain bound]
Under the setup above,
\[
  \KL{P}{Q} \leq \sum_{i<n}\varepsilon_i.
\]
\end{lemma}

\begin{proof}[Proof plan]
For a trace $x=(x_0,\ldots,x_{n-1})$, write
$x_{<i}=(x_0,\ldots,x_{i-1})$.  On the support of $P$, absolute continuity
allows the usual chain factorization
\[
  \log\frac{P(x)}{Q(x)}
  =
  \sum_{i<n}
    \log
      \frac{P_i(x_i \mid x_{<i})}
           {Q_i(x_i \mid x_{<i})}.
\]
Multiplying by $P(x)$ and summing over $x$ gives
\[
\begin{aligned}
  \KL{P}{Q}
  &=
  \sum_{i<n}
  \sum_{x\in A^n}
    P(x)
    \log
      \frac{P_i(x_i \mid x_{<i})}
           {Q_i(x_i \mid x_{<i})}.
\end{aligned}
\]
For a fixed coordinate $i$, group traces by their prefix $a=x_{<i}$ and
current coordinate value $b=x_i$.  Summing out the suffix yields
\[
  \sum_{x\in A^n}
    P(x)
    \log
      \frac{P_i(x_i \mid x_{<i})}
           {Q_i(x_i \mid x_{<i})}
  =
  \sum_{a\in A^i} P(x_{<i}=a)\,
    \KL{P_i(\cdot\mid a)}{Q_i(\cdot\mid a)}.
\]
By the coordinate hypothesis this is at most
\[
  \sum_{a\in A^i} P(x_{<i}=a)\,\varepsilon_i
  =
  \varepsilon_i,
\]
because prefix masses sum to $1$.  Summing over coordinates gives
\[
  \KL{P}{Q}
  \leq
  \sum_{i<n}\varepsilon_i.
\]
\end{proof}

\begin{lemma}[Bounded coordinate KL implies summability]
Under the setup above, with finite conditional KL divergences,
\[
  \sum_{x \in A^n}
    \left|P(x)\log\frac{P(x)}{Q(x)}\right|
\]
is finite.  More concretely, for every finite set $J \subseteq A^n$,
\[
  \sum_{x \in J}
    \left(P(x)\log\frac{P(x)}{Q(x)}\right)^+
  \leq
  \sum_{i < n} \varepsilon_i + n.
\]
\end{lemma}

\begin{proof}[Proof plan]
First observe that the negative part is already harmless once
$P \ll Q$.  Pointwise,
\[
  P(x)\log\frac{P(x)}{Q(x)} \geq -Q(x),
\]
so the negative part is bounded by $Q(x)$ and is therefore summable.  Thus it
is enough to bound finite sums of the positive part.

For a trace $x=(x_0,\ldots,x_{n-1})$, write
$x_{<i}=(x_0,\ldots,x_{i-1})$.  Reuse the same chain factorization as in the
value-level proof:
\[
  \log\frac{P(x)}{Q(x)}
  =
  \sum_{i<n}
    \log
      \frac{P_i(x_i \mid x_{<i})}
           {Q_i(x_i \mid x_{<i})}.
\]
Hence, using $(u+v)^+ \leq u^+ + v^+$,
\[
  \left(P(x)\log\frac{P(x)}{Q(x)}\right)^+
  \leq
  \sum_{i<n}
    P(x)
    \left(
      \log
      \frac{P_i(x_i \mid x_{<i})}
           {Q_i(x_i \mid x_{<i})}
    \right)^+.
\]
Fix a coordinate $i$ and sum the $i$th term over a finite set
$J \subseteq A^n$.  Group the terms by prefix $a \in A^i$ and current
coordinate value $b \in A$.  The remaining suffix mass inside $J$ is at most
the full conditional suffix mass, so
\[
\begin{aligned}
  &\sum_{x \in J}
    P(x)
    \left(
      \log
      \frac{P_i(x_i \mid x_{<i})}
           {Q_i(x_i \mid x_{<i})}
    \right)^+
  \\
  &\qquad\leq
  \sum_{a \in A^i} P(x_{<i}=a)
  \sum_{b \in A}
    P_i(b \mid a)
    \left(
      \log\frac{P_i(b \mid a)}{Q_i(b \mid a)}
    \right)^+.
\end{aligned}
\]
For each prefix $a$, let
\[
  h_{i,a}(b)
  =
  P_i(b \mid a)\log\frac{P_i(b \mid a)}{Q_i(b \mid a)}.
\]
Since the conditional KL is finite,
\[
  \sum_b h_{i,a}(b)^+
  =
  \KL{P_i(\cdot \mid a)}{Q_i(\cdot \mid a)}
  +
  \sum_b h_{i,a}(b)^-.
\]
The same lower bound $h_{i,a}(b)\geq -Q_i(b\mid a)$ gives
\[
  \sum_b h_{i,a}(b)^-
  \leq
  \sum_b Q_i(b \mid a)
  \leq 1.
\]
Therefore
\[
  \sum_b h_{i,a}(b)^+
  \leq
  \varepsilon_i + 1.
\]
Averaging over prefixes with weights $P(x_{<i}=a)$, whose total mass is $1$,
gives an $i$th-coordinate contribution at most $\varepsilon_i+1$.  Summing over
all $i<n$ gives
\[
  \sum_{x \in J}
    \left(P(x)\log\frac{P(x)}{Q(x)}\right)^+
  \leq
  \sum_{i<n}(\varepsilon_i+1)
  =
  \sum_{i<n}\varepsilon_i+n.
\]
This is a uniform bound over finite $J$, which is exactly the summability
criterion for the positive part.  Together with the negative-part bound, it
implies absolute summability of the KL integrand.
\end{proof}

\paragraph{Rocq translation notes.}
The value-level proof corresponds to \texttt{iterated\_kl\_chain\_bound}.  The
summability proof corresponds to
\texttt{finite\_sum\_fpos\_kl\_from\_coordinate\_bounded\_kl} and then
\texttt{summable\_kl\_from\_coordinate\_bounded\_kl}.  The shared engineering
work is the chain decomposition/grouping step by prefix and coordinate value.
For the value theorem, this grouping produces an expectation of conditional
KLs.  For the summability theorem, the same grouping is applied to finite
positive-part sums and pays one negative-part slack per coordinate.  The
already available lemmas \texttt{summable\_fneg\_kl\_integrand} and
\texttt{summable\_from\_fpos\_fneg} handle the final passage from positive-part
summability to KL-integrand summability.
